\originalchapter{8.321研究生作业七至十二}
以下题面据官方PDF中文转述，解答为\textbf{AI独立解答}；研究生题旁的商业教材改编信用不移除。磁场计算用SI制，原件高斯制的$qB/c$在此换成$qB$。

\originalsection{作业7：磁平移、色散与中性两体}
\begin{exercise}
PS7 Q1：对二维电荷$e$的$\vct B=B\hat z$取对称规范$A_x=-By/2,A_y=Bx/2$。(a)核磁场；(b)联系二维谐振子求不同能量；(c)用$z=x+iy$求最低Landau态；(d)可选题：半径$R$圆盘的最低能级简并数。
\end{exercise}
\begin{solution}
$\partial_xA_y-\partial_yA_x=B$。展开$H=p^2/(2m)+m\omega_c^2r^2/8-eBL_z/(2m)$，其中$\omega_c=|eB|/m$。二维频率$\omega_c/2$谐振子的谱$\hbar\omega_c(2n_r+|l|+1)/2-\hbar eBl/(2m)$，等价$E_n=\hbar\omega_c(n+1/2)$。假定$eB>0$，最低层是$n_r=0,l\ge0$，基函数
$\psi_l=(2\pi2^ll!\ell_B^2)^{-1/2}(z/\ell_B)^l e^{-r^2/(4\ell_B^2)}$、$\ell_B^2=\hbar/(eB)$；正规最低层一般为全纯函数乘该Gaussian，$eB<0$时改反全纯。径向密度峰约$r^2=2l\ell_B^2$，圆盘内$l\lesssim R^2/(2\ell_B^2)$，简并数约$R^2/(2\ell_B^2)=|eB|\pi R^2/h$，边缘修正不计；与Landau规范计数一致。
\end{solution}
\begin{exercise}
PS7 Q2：色散改为$\epsilon(\Pi)=\Pi^2/(2m)+\lambda(\Pi^2)^2$、$\lambda>0$，证明各能级中心简并不变。Q3：三维恒$B\hat z$加$E\hat x$，用Landau规范证明$p_y,p_z$守恒，求精确能谱及群速度。
\end{exercise}
\begin{solution}
Q2导引中心$R_i$与$\Pi_j$对易，因此与任意该色散对易；回旋本征态没有改变。$\Pi^2=2\hbar|eB|(n+1/2)$，其能量映射严格单调，故没有另生层间简并，中心计数仍为磁通量子数。

Q3取$H=[p_x^2+(p_y-eBx)^2+p_z^2]/(2m)-eEx$，没有$y,z$，故两正则动量守恒。设$x_0=p_y/(eB)$，配方中心为$x_c=x_0+mE/(eB^2)$，得
\[
 E_{n,p_y,p_z}=\hbar\omega_c(n+1/2)+\frac{p_z^2}{2m}
 -\frac EBp_y-\frac{mE^2}{2B^2}.
\]
原Landau层随$p_y$倾斜，均匀电场全空间模型没有全局基态。$v_z=\partial E/\partial p_z=p_z/m$、$v_y=-E/B$，正是$\vct E\times\vct B/B^2$漂移。常项的负号来自平移后的势能，不等于机械动能负。
\end{solution}
\begin{exercise}
PS7 Q4：两等质量$m$粒子电荷$\pm q$、均匀$B$、相互作用$V=\kappa|\vct x_1-\vct x_2|^2/2$。取$R=(x_1+x_2)/2,x=x_1-x_2$，$Q=\Pi_1+\Pi_2-q\vct x\times\vct B$、$p=(\Pi_1-\Pi_2)/2$。核正则对易、$Q$守恒及经典对应；固定$Q$改写$H$，求$Q=0$谱。
\end{exercise}
\begin{solution}
对称规范下，若$P=p_1+p_2$、$p_r=(p_1-p_2)/2$，则$Q=P+q\vct B\times x/2$、$p=p_r-q\vct B\times R/2$。用$[R_i,P_j]=[x_i,p_{rj}]=i\hbar\delta_{ij}$，其他交叉为0，立即得$[R_i,Q_j]=[x_i,p_j]=i\hbar\delta_{ij}$，并有$[Q_i,p_j]=[Q_i,x_j]=[Q_i,Q_j]=[p_i,p_j]=[R_i,p_j]=0$。交叉$[Q,p]$中的两磁项互相抵消。

解$\Pi_1+\Pi_2=Q+q x\times B$、$\Pi_1-\Pi_2=2p$给$H=(Q+qx\times B)^2/(4m)+p^2/m+V(x)$。它与$Q$交换；经典式$\frac{\dd}{\dd t}(\Pi_1+\Pi_2)=q(v_1-v_2)\times B$使同一$Q$恒定，内力抵消。$Q=0$相对二维谐振子质量$\mu=m/2$、频率$\Omega=\sqrt{(qB/m)^2+2\kappa/m}$，故$E=\hbar\Omega(n_x+n_y+1)$，每个总激发数$N$有$N+1$简并。固定$Q$是广义本征扇区，中心坐标平面波不当作全空间归一束缚态。
\end{solution}

\originalsection{作业8：偏迹、净化与熵}
\begin{exercise}
PS8 Q1：三自旋态$N(\ket{\uparrow\downarrow\uparrow}+\ket{\downarrow\uparrow\downarrow}+\ket{\uparrow\uparrow\uparrow})$，求$N$、各自旋$S_z$正读数概率、$A$密度矩阵，并由它求$S_z,S_x$概率。Q3：$\ket\psi=(\ket{0,+}+\ket{1,0})/\sqrt2$，求Schmidt分解和两边密度矩阵。
\end{exercise}
\begin{solution}
Q1三积态正交，$N=1/\sqrt3$。三个自旋的$+z$概率都为$2/3$；偏迹$BC$使$A$交叉项全消失，$\rho_A=\diag(2/3,1/3)$。故$S_z$两概率$2/3,1/3$，$S_x$两概率均$1/2$。

Q3在$z$基的系数矩阵为$C=\frac12\begin{pmatrix}1&1\\\sqrt2&0\end{pmatrix}$，所以
$\rho_A=CC^\dagger=\begin{pmatrix}1/2&1/(2\sqrt2)\\1/(2\sqrt2)&1/2\end{pmatrix}$，
$\rho_B=C^TC^*=\begin{pmatrix}3/4&1/4\\1/4&1/4\end{pmatrix}$。
非零本征值$\lambda_\pm=(1\pm1/\sqrt2)/2$；$u_\pm=(\ket0\pm\ket1)/\sqrt2$，$v_+=\cos(\pi/8)\ket0+\sin(\pi/8)\ket1$、$v_-=-\sin(\pi/8)\ket0+\cos(\pi/8)\ket1$。于是$\ket\psi=\sqrt{\lambda_+}u_+v_++\sqrt{\lambda_-}u_-v_-$，展开即可复核原系数。
\end{solution}
\begin{exercise}
PS8 Q4：证明同一$\rho_A$在同一辅助系统$B$上的任意两种纯态净化，仅差$B$上的幺正。
\end{exercise}
\begin{solution}
先取本题隐含的有限维系统。选$\rho_A$本征基$\ket j$，任意净化可写$\sum_j\sqrt{\lambda_j}\ket j\ket{b_j}$；偏迹条件要求非零$\lambda_j$的$\ket{b_j}$正交归一。另一净化同样为$\sum\sqrt{\lambda_j}\ket j\ket{b'_j}$。在这两个等维支撑子空间定义$U_Bb_j=b'_j$，再选两侧正交补基逐一对应，扩成整个$B$的幺正，便得$\psi'=(I\otimes U_B)\psi$。若辅助空间本来不同，只能先说明等距映射或补足维数，不能无条件称同一幺正。无限维同一辅助空间也须两支撑的正交补等维才可作幺正延拓；否则只保证支撑间等距映射，不能无限外推原题断言。
\end{solution}
\begin{exercise}
PS8 Q5：对$S(\rho)=-\Tr\rho\ln\rho$证明(a)幺正不变；(b)秩$D$时$S\le\ln D$并给等号；(c)密度矩阵凸混合仍合法且熵凹；(d)$S_{AB}\le S_A+S_B$及等号条件。
\end{exercise}
\begin{solution}
(a)幺正不改谱，熵是$-\sum r_i\ln r_i$，约定$0\ln0=0$。(b)对均匀$u_i=1/D$，标量相对熵$\sum r_i\ln(r_i/u_i)\ge0$给$S\le\ln D$，等号仅所有非零$r_i=1/D$。

为证明(c)(d)，先给量子相对熵的非负证明。将$\rho=\sum r_i\ket i\bra i$、$\sigma=\sum s_j\ket{v_j}\bra{v_j}$，假设$\operatorname{supp}\rho\subset\operatorname{supp}\sigma$。对每$i$，$w_{ij}=|\braket i{v_j}|^2$和为1，$\ln$凹性给$\sum_jw_{ij}\ln s_j\le\ln\sigma_{ii}$，所以
$D(\rho\|\sigma)\ge\sum_i r_i\ln(r_i/\sigma_{ii})\ge0$。
后一式用$\ln x\le x-1$，$\sum\sigma_{ii}=1$；不满足支撑条件时$D=+\infty$。等号条件合起来是$\rho=\sigma$。

(c)$\rho_\lambda=\lambda\rho_1+(1-\lambda)\rho_2$对$0\le\lambda\le1$保Hermitian、正性、迹1。恒等式$S(\rho_\lambda)-\lambda S(\rho_1)-(1-\lambda)S(\rho_2)=\lambda D(\rho_1\|\rho_\lambda)+(1-\lambda)D(\rho_2\|\rho_\lambda)$给凹性。(d)利用张量积对数$\ln(\rho_A\otimes\rho_B)=\ln\rho_A\otimes I+I\otimes\ln\rho_B$及偏迹，得$D(\rho_{AB}\|\rho_A\otimes\rho_B)=S_A+S_B-S_{AB}\ge0$，等号当且仅当$\rho_{AB}=\rho_A\otimes\rho_B$。零本征值在支撑上取极限，不能直接对零取普通对数。
\end{solution}

\originalsection{作业9：Wigner的限定版本与角动量代数}
\begin{exercise}
PS9 Q1：证明保持所有测量转移概率的线性或反线性对称由幺正/反幺正实现。Q2：反演与平移为何不交换，其对自由谱有何意义？Q3：证明$[J^2,J_i]=0$、一般向量算符$V^2$与$J_i$对易，并由$x,p$基本对易核$L=x\times p$的代数。
\end{exercise}
\begin{solution}
Q1若线性$U$按原题保持规范化态的全部内积模长并将单位态映成单位态，对任意$v$齐次性给$\|Uv\|=\|v\|$；极化恒等式从四个范数$\|u\pm v\|$、$\|u\pm iv\|$重建复内积，故$U^\dagger U=I$。再由对称的可逆/满射性得$UU^\dagger=I$。反线性时选定正交基复共轭$K$，$UK$线性且等距，因此$U=(UK)K$反幺正。原件将反线性也称“unitary”措辞不精确；没有假设线性/反线性的完整射线Wigner定理比此题限定版本更强，此处不冒称已从转移概率推导其全部前提。若只把输出归一后比较概率，还须先排除任意整体幅度尺度，才讨论幺正代表。

Q2$PT(a)P^{-1}=T(-a)$，一般不与$T(a)$交换。自由$H=p^2/(2m)$同时有两对称，$p\ne0$态与$-p$同能，反演把其交换；$p=0$不必成对。Q3$[J_i,\sum_jJ_jJ_j]=i\hbar\sum_{jk}\epsilon_{ijk}(J_kJ_j+J_jJ_k)=0$，对称乘积与反对称指标收缩为零。同理向量算符满足$[J_i,V_j]=i\hbar\epsilon_{ijk}V_k$时，$[J_i,V^2]=0$，不需$V_j$互相交换。轨道计算先得$[x_ap_b,x_cp_d]=i\hbar(\delta_{ad}x_cp_b-\delta_{bc}x_ap_d)$，再双收缩$\epsilon_{iab}\epsilon_{jcd}$得$[L_i,L_j]=i\hbar\epsilon_{ijk}L_k$。
\end{solution}

\originalsection{作业10：有限旋转与角动量合成}
\begin{exercise}
PS10 Q1：核自旋$1/2$对易式，写三轴有限旋转矩阵，将$\ket{+x}$绕$z$转$\phi$后求其正自旋方向。Q2：证明$R_{ij}=\Tr(\sigma_jU^\dagger\sigma_iU)/2$对$U\in SU(2)$给SO(3)，$U,-U$给相同旋转。
\end{exercise}
\begin{solution}
$[\sigma_i,\sigma_j]=2i\epsilon_{ijk}\sigma_k$给$[S_i,S_j]=i\hbar\epsilon_{ijk}S_k$；原Q1(a)漏$i$，此处校正。令$c=\cos(\theta/2),s=\sin(\theta/2)$，三轴矩阵为
$D_x=\begin{pmatrix}c&-is\\-is&c\end{pmatrix}$、$D_y=\begin{pmatrix}c&-s\\s&c\end{pmatrix}$、$D_z=\diag(e^{-i\theta/2},e^{i\theta/2})$；原轴列表重复$x$，按三坐标轴补$z$。$D_z(\phi)\ket{+x}$的Bloch向量$(\cos\phi,\sin\phi,0)$。

$U^\dagger\sigma_iU$为无迹Hermitian，可展开$R_{ij}\sigma_j$，系数实。迹内积保持使$R_{ij}R_{kj}=\delta_{ik}$，故$R$正交。SU(2)连通且$U=I$时行列式1，不能连续跳到$-1$，所以$R\in SO(3)$。表达式中的两负号抵消，$U,-U$映成相同$R$。
\end{solution}
\begin{exercise}
PS10 Q3：自旋1绕$x$转动，证明$D_x=I-i\sin\theta S_x/\hbar+(\cos\theta-1)(S_x/\hbar)^2$；初$l=1,m=1$分别绕$x,z$转后测$m=1$概率。Q4：$l$与电子自旋$1/2$合成，求允许$j$及全部CG系数。
\end{exercise}
\begin{solution}
$M=S_x/\hbar$谱$-1,0,1$，所以$M^3=M$；指数奇次求和为$-iM\sin\theta$、正偶次为$(\cos\theta-1)M^2$。在$S_z$基$\langle1|M|1\rangle=0$、$\langle1|M^2|1\rangle=1/2$，留在$m=1$的幅度$(1+\cos\theta)/2$，概率$\cos^4(\theta/2)$；绕$z$只有相位，概率1。

$j=l+1/2$或$l-1/2$（$l=0$只有前者）。固定总$m$在两积态$\ket{m-1/2,+1/2}$、$\ket{m+1/2,-1/2}$中，取Condon--Shortley相位：
\[
\begin{aligned}
\ket{l+\tfrac12,m}&=\sqrt{\frac{l+m+1/2}{2l+1}}\ket{m-1/2,+1/2}
+\sqrt{\frac{l-m+1/2}{2l+1}}\ket{m+1/2,-1/2},\\
\ket{l-\tfrac12,m}&=-\sqrt{\frac{l-m+1/2}{2l+1}}\ket{m-1/2,+1/2}
+\sqrt{\frac{l+m+1/2}{2l+1}}\ket{m+1/2,-1/2}.
\end{aligned}
\]
从最大权态$\ket{l,l}\ket{+1/2}$逐次作用$J_-=L_-+S_-$给第一式，第二式是同$m$二维空间的归一正交补；只有合法积态和$m$范围使用。相位可整体改变，不改测量概率。
\end{solution}

\originalsection{作业11：非谐振子与简并微扰}
\begin{exercise}
PS11 Q2（Sakurai5.1改编）：谐振子加$bx$，求基态最低非零能移、精确解；再分别加$\lambda x^3$或$\lambda x^4$求各能级最低非零能移。
\end{exercise}
\begin{solution}
设$d=\sqrt{\hbar/(2m\omega)}$，$x=d(a+a^\dagger)$。线性场只连基态到1态，$\Delta E_0^{(2)}=-b^2d^2/(\hbar\omega)=-b^2/(2m\omega^2)$。配方$V=m\omega^2(x+b/(m\omega^2))^2/2-b^2/(2m\omega^2)$，所有态有同一精确能移。

三次项一阶因宇称为0；$(a+a^\dagger)^3$作用在$n$态的$n-3,n-1,n+1,n+3$系数分别为$\sqrt{n(n-1)(n-2)},3n\sqrt n,3(n+1)\sqrt{n+1},\sqrt{(n+1)(n+2)(n+3)}$。代入二阶能量和得
$\Delta E_n=-\lambda^2\hbar^2(30n^2+30n+11)/(8m^3\omega^4)$。
四次项一阶矩元$\langle(a+a^\dagger)^4\rangle=6n^2+6n+3$，故$\Delta E_n=3\lambda d^4(2n^2+2n+1)$。纯三次势全直线上无下界，所给是弱非谐局部模型的形式微扰，不能称真实全局束缚谱；正四次势则稳定。
\end{solution}
\begin{exercise}
PS11 Q4：$H_0=-J\vct S_1\cdot\vct S_2$、$J>0$，自旋均$1/2$；微扰$H_1=-BS_{1z}$只作用第一自旋，求谱及小$B$的首阶变化。
\end{exercise}
\begin{solution}
总自旋$t=1$三重态能$-J\hbar^2/4$、$s=0$单态能$3J\hbar^2/4$。三重空间中$PS_{1z}P=J_z/2$，故$m=1,0,-1$一阶能移$-B\hbar m/2$；单态一阶为0。这里真正小量条件为$|B|\hbar\ll J\hbar^2$，原写$B\ll J$隐含$\hbar=1$。$m=0$三重态与单态耦合$\langle s|H_1|t_0\rangle=-B\hbar/2$，若要求这些首阶零的态最低非零结果，则分别为$\mp B^2/(4J)$。其精确二维谱为$J\hbar^2/4\pm\frac12\sqrt{J^2\hbar^4+B^2\hbar^2}$，展开验证上述结果。
\end{solution}

\originalsection{未取得题面与练习年代}
PS8 Q2=Sakurai3.10；PS9 Q4=3.5、Q5=3.14；PS10 Q5=4.2；PS11 Q1=4.11、Q3=5.2、Q5=5.9、Q6=5.20；PS12 Q1--5=5.25、5.28、5.30、5.35、5.37。共13道纯教材引用。PS12原件正文标Fall 2016、For practice--do not submit，公开资源页在Fall 2017课程；保留该年代及性质，不按2017计分作业伪称答案已完成。本研究生8.321作业中共20道教材引用（前组7道、本组13道），与旧本科卷10道缺项分别记录。
